% §2 Preliminaries
%
% Notation, fermionic algebra, Pauli algebra, Majorana operators.
% Kept self-contained for tutorial accessibility.

We establish the notation and collect the key definitions used
throughout.

\subsection{Fermionic operators and canonical anticommutation relations}

Consider $n$ fermionic modes with creation operators $\ad{j}$ and
annihilation operators $\an{j}$ for $j \in \{0, \ldots, n{-}1\}$.
These satisfy the \emph{canonical anticommutation relations} (CAR):
\begin{align}
  \acomm{\an{i}}{\ad{j}} &= \delta_{ij}, \label{eq:car1} \\
  \acomm{\an{i}}{\an{j}} &= 0, \label{eq:car2} \\
  \acomm{\ad{i}}{\ad{j}} &= 0. \label{eq:car3}
\end{align}
These relations encode the Pauli exclusion principle: no two fermions
can occupy the same mode.  The \emph{number operator}
$\nhat{j} = \ad{j}\an{j}$ has eigenvalues 0 and~1, and the Fock space
basis states are tensor products $\ket{f_{n-1} \cdots f_1 f_0}$ with
$f_j \in \{0, 1\}$.

\subsection{Pauli algebra}
\label{sec:pauli-algebra}

A qubit register of $n$ qubits admits the Pauli group $\mathcal{P}_n$
consisting of $n$-fold tensor products of single-qubit Pauli operators
$\{I, X, Y, Z\}$ with phases in $\{+1, -1, +i, -i\}$.  The
single-qubit multiplication rules are:
\begin{equation}
  XY = iZ, \quad YZ = iX, \quad ZX = iY,
  \label{eq:pauli-products}
\end{equation}
with $\sigma^2 = I$ for all $\sigma \in \{X, Y, Z\}$ and
$\sigma_a \sigma_b = -\sigma_b \sigma_a$ for $a \neq b$.

Multi-qubit Pauli strings multiply componentwise:
\begin{equation}
  (\sigma_{a_0} \otimes \cdots \otimes \sigma_{a_{n-1}})
  (\sigma_{b_0} \otimes \cdots \otimes \sigma_{b_{n-1}})
  = \phi \cdot
  (\sigma_{a_0}\sigma_{b_0}) \otimes \cdots \otimes
  (\sigma_{a_{n-1}}\sigma_{b_{n-1}}),
  \label{eq:pauli-mult}
\end{equation}
where $\phi = \prod_{k=0}^{n-1} \phi_k$ is the accumulated phase from
each single-qubit multiplication.

\begin{definition}[Pauli weight]
\label{def:pauli-weight}
The \emph{Pauli weight} of a string $\sigma \in \mathcal{P}_n$ is the
number of non-identity tensor factors.  For an encoded ladder operator
expressed as a sum of Pauli strings, the weight is the maximum over all
terms.
\end{definition}

Pauli weight directly determines the number of CNOT gates required to
measure the corresponding term in a variational quantum eigensolver
(VQE) circuit, making it the primary figure of merit for encoding
quality.

\subsection{Majorana operators}
\label{sec:majorana}

\begin{definition}[Majorana operators]
\label{def:majorana}
For each mode $j$, the Hermitian \emph{Majorana operators} are:
\begin{align}
  c_j &= \ad{j} + \an{j}, \label{eq:majorana-c} \\
  d_j &= i(\ad{j} - \an{j}). \label{eq:majorana-d}
\end{align}
\end{definition}

These satisfy the Majorana anticommutation relations:
\begin{equation}
  \{c_i, c_j\} = 2\delta_{ij}, \quad
  \{d_i, d_j\} = 2\delta_{ij}, \quad
  \{c_i, d_j\} = 0.
  \label{eq:majorana-car}
\end{equation}

Inverting~\eqref{eq:majorana-c}--\eqref{eq:majorana-d}:
\begin{equation}
  \ad{j} = \tfrac{1}{2}(c_j - id_j), \qquad
  \an{j} = \tfrac{1}{2}(c_j + id_j).
  \label{eq:ladder-from-majorana}
\end{equation}

The Majorana decomposition is the key bridge between fermions and
qubits: every encoding reduces to specifying how to represent $c_j$
and $d_j$ as Pauli strings.  If the Pauli representations satisfy
the Majorana anticommutation relations~\eqref{eq:majorana-car}, the
induced ladder operators automatically satisfy the
CAR~\eqref{eq:car1}--\eqref{eq:car3}.

\begin{definition}[Fermion-to-qubit encoding]
\label{def:encoding}
A \emph{fermion-to-qubit encoding} for $n$ modes on $n$ qubits is a
pair of maps $c_j \mapsto S_j^{(c)}$ and $d_j \mapsto S_j^{(d)}$, where
each $S_j^{(c)}$ and $S_j^{(d)}$ is a Pauli string in
$\mathcal{P}_n$, such that the induced ladder operators
\begin{equation}
  \ad{j} = \tfrac{1}{2}(S_j^{(c)} - i S_j^{(d)}), \qquad
  \an{j} = \tfrac{1}{2}(S_j^{(c)} + i S_j^{(d)})
  \label{eq:encoded-ladder}
\end{equation}
satisfy the CAR~\eqref{eq:car1}--\eqref{eq:car3}.
\end{definition}

\subsection{The electronic Hamiltonian}
\label{sec:hamiltonian}

The second-quantized electronic Hamiltonian is:
\begin{equation}
  H = \sum_{pq} h_{pq}\, \ad{p}\an{q}
    + \frac{1}{2} \sum_{pqrs} g_{pqrs}\, \ad{p}\ad{q}\an{s}\an{r},
  \label{eq:hamiltonian}
\end{equation}
where $h_{pq}$ are one-electron integrals and $g_{pqrs}$ are
two-electron integrals from quantum chemistry
computations~\cite{szabo1996, helgaker2000}.  An encoding maps each
$\ad{j}$ and $\an{j}$ to a two-term Pauli expression
via~\eqref{eq:encoded-ladder}, and the products are expanded using the
Pauli multiplication rule~\eqref{eq:pauli-mult}.  The result is a
Hamiltonian expressed entirely in terms of Pauli strings:
\begin{equation}
  H = \sum_{\alpha} h_\alpha\, \sigma_\alpha,
  \label{eq:pauli-hamiltonian}
\end{equation}
where $\sigma_\alpha \in \mathcal{P}_n$ and $h_\alpha \in \mathbb{C}$.

The choice of encoding affects neither the eigenspectrum nor the
physics, only the measurement cost: each distinct Pauli string
requires a separate measurement basis, and the number of such terms
and their weights determine the overall simulation overhead.
